\section{Биохимия задачи}
\label{sec:bio}

Раздел для тех, кто пришёл со стороны машинного обучения: что физически стоит за каждой колонкой
в данных. Без этого половина решений в модели выглядит произволом.

\subsection{Что такое цитохромы P450 и почему именно эти четыре}

Цитохромы P450 --- надсемейство гем-тиолатных монооксигеназ. Гем --- порфириновое кольцо с железом
в центре, то же самое, что в гемоглобине; «тиолатные» означает, что пятое координационное место железа
занимает атом серы остатка цистеина. Название историческое: восстановленный комплекс фермента с угарным
газом даёт полосу поглощения при 450 нм.

Каталитический цикл в общем виде: фермент связывает субстрат, железо восстанавливается, присоединяет
кислород, кислород расщепляется с образованием высокореакционного оксоферрильного интермедиата
(Compound~I), который вырывает водород у субстрата и вставляет туда гидроксил. Суммарно
\[
\text{RH} + \text{O}_2 + \text{NADPH} + \text{H}^+ \;\longrightarrow\;
\text{ROH} + \text{H}_2\text{O} + \text{NADP}^+ .
\]
Электроны поставляет НАДФН через флавопротеин цитохром-P450-редуктазу. Для эксперимента отсюда следует вот что: без НАДФН фермент не работает, так что катализ можно включать
и выключать в пробирке. На этом приёме построено измерение TDI.

Смысл реакции --- сделать липофильную чужеродную молекулу более полярной, чтобы её можно было
конъюгировать (вторая фаза метаболизма) и вывести. Печёночные P450 --- главный барьер между
проглоченным веществом и системным кровотоком.

Изоформ у человека около 57, но лекарственный метаболизм сосредоточен в немногих. CYP3A4 отвечает
примерно за половину всех оборотов, CYP2D6, CYP2C9, CYP2C19 и CYP1A2 --- за большую часть остального.
Четвёрка в соревновании покрывает подавляющее большинство клинически значимых взаимодействий.

\subsection{Почему ингибирование --- проблема}

Если препарат A ингибирует фермент, который выводит препарат B, концентрация B в крови растёт ---
иногда в разы. Для препарата с узким терапевтическим окном это прямая дорога к токсичности, вплоть
до снятия с рынка уже зарегистрированного средства, безопасного само по себе.

Поэтому ингибирование P450 проверяют рано, ещё на стадии выбора соединения-кандидата. Предсказание
по структуре экономит эти пробы, и цена ошибки здесь измеряется в снятых с разработки молекулах.

\subsection{Что такое IC$_{50}$ и pIC$_{50}$}

В пробу вносят фермент, зонд-субстрат с известной скоростью превращения и испытуемое соединение
в нескольких концентрациях. Измеряют, насколько упала скорость превращения зонда. IC$_{50}$ ---
концентрация испытуемого соединения, при которой скорость падает вдвое. $\pic = -\log_{10}$IC$_{50}$
в молях на литр: $\pic = 6$ соответствует одному микромолю, $\pic = 4$ --- ста микромолям. Больше $\pic$ ---
сильнее ингибитор.

Важная оговорка: IC$_{50}$ --- не константа связывания. Она зависит от концентрации зонд-субстрата
и от механизма (конкурентный ингибитор вытесняется субстратом, неконкурентный --- нет); связь
с константой ингибирования $K_i$ даёт соотношение Ченга--Прусоффа. Сравнивать IC$_{50}$ между
лабораториями и пробами напрямую нельзя, внутри одного набора --- можно, и именно это мы и делаем.

\subsection{Как читают результат и откуда берутся артефакты}

Зонд-субстрат подбирают так, чтобы его превращение было легко измерить. Способов два.

\emph{Флуоресцентный}: зонд после окисления даёт флуоресцирующий продукт, и прибор просто меряет
свечение. Дёшево и быстро, поэтому так сделаны 1A2, 2C9 и 3A4 в этом наборе.

\emph{Масс-спектрометрический}: продукт детектируют по массе. Дороже, но специфичнее; так сделан 2D6.

Отсюда конкретный артефакт. Если испытуемое соединение само флуоресцирует в том же диапазоне или,
наоборот, тушит флуоресценцию продукта, прибор покажет изменение сигнала, которого химически не было.
Ложное «ингибирование» появится в трёх флуоресцентных эндпойнтах и не появится в масс-спектрометрическом.
Поэтому в модели заведён отдельный канал артефактов пробы (раздел~\ref{sec:artifact}): это конфаундер с известной
структурой, а не шум.

\subsection{Что такое time-dependent ингибирование}
\label{sec:tdi}

Обычное ингибирование обратимо: молекула садится в карман, мешает субстрату, отпускает. Уберите
ингибитор --- активность вернётся.

Time-dependent ингибирование, оно же механизм-зависимая инактивация, устроено иначе. Фермент сам
превращает соединение в реакционноспособный продукт, и этот продукт остаётся --- ковалентно
модифицирует гем или белок, либо образует прочный комплекс с железом (metabolite-intermediate complex).
Фермент выведен из строя необратимо. Отличительные признаки: процесс требует катализа (без НАДФН
не идёт), развивается во времени, насыщается по концентрации, и активность не восстанавливается
отмывкой.

Клинически это гораздо хуже обратимого ингибирования. Обратимый ингибитор перестаёт мешать, как только
выводится из организма. Инактивированный фермент восстанавливается только синтезом нового белка,
а на это уходят сутки-двое. Поэтому эффект накапливается при повторных приёмах и держится долго после
отмены.

\paragraph{Как это измеряют.} Ставят два плеча одной пробы.

\emph{Прямое плечо}: соединение и зонд-субстрат вносят одновременно, сразу меряют IC$_{50}$. Видно
обратимое ингибирование.

\emph{Плечо с преинкубацией}: соединение сначала выдерживают с ферментом и НАДФН, дав катализу
поработать, и только потом вносят зонд и меряют IC$_{50}$. Если соединение --- механизм-зависимый
инактиватор, за время преинкубации часть фермента выведена из строя, и измеренная IC$_{50}$ окажется
ниже: соединение выглядит потентнее, чем было.

Разница между плечами и есть сигнал. В логарифмической шкале это сдвиг
\[
\Delta = \pi_{\text{после преинкубации}} - \pi_{\text{прямое}} \;\ge\; 0 ,
\]
и неотрицательность здесь не соглашение, а физика: преинкубация может только добавить инактивации,
но не отменить обратимое связывание. Двукратный сдвиг IC$_{50}$ --- то есть $\Delta > \log_{10}2 \approx 0.301$
--- общепринятый порог, за которым соединение помечают как TDI.

Отсюда становится понятной вторая ветвь правила разметки из раздела~\ref{sec:decision}. Соединение может
быть настолько слабым обратимым ингибитором, что в прямом плече его IC$_{50}$ вообще не определяется
как значимая ($\pi \le 4$, то есть слабее ста микромолей). Но если после преинкубации оно перешло порог
активности с тем же двукратным запасом ($\pi + \Delta > 4 + \log_{10}2$), это тоже TDI --- просто
инактивация оказалась единственным его механизмом. Правило и состоит из этих двух случаев.

\subsection{Что такое Emax и положительный контроль}

Кривая доза-эффект выходит на плато. Высота плато --- Emax, предельная доля подавленной активности.
Полный ингибитор обнуляет активность, частичный оставляет какую-то долю сколь угодно высокой
концентрации. Причины частичности бывают разные: соединение садится не в каталитический карман,
а в аллостерический; или у фермента несколько сайтов связывания субстрата, как у CYP3A4.

Измеряют Emax не в абсолютных единицах, а относительно положительного контроля --- известного полного
ингибитора этого фермента, который ставят в ту же пробу. Отсюда шкала в данных: полный ингибитор даёт
примерно $-1$, потому что подавил столько же, сколько контроль. Такая нормировка снимает разброс между
планшетами и партиями фермента.

\subsection{Чем четыре фермента различаются}
\label{sec:isoforms}

Отсюда понятно, почему модель устроена как многозадачная с явной матрицей связей, а не как одна голова.

\emph{CYP3A4} --- самый крупный и самый неразборчивый. Активный карман большой и податливый, вмещает
молекулы вплоть до циклических пептидов; в нём помещается больше одного лиганда одновременно, откуда
кооперативность и не всегда сигмоидные кривые. Узнавание в основном гидрофобное, поэтому предсказатель
«крупная жирная молекула --- вероятный ингибитор» работает удивительно хорошо. Отсюда и лучшая
предсказуемость 3A4 в наших измерениях.

\emph{CYP2C9} --- карман поменьше, на входе положительно заряженный остаток аргинина, притягивающий
анионы. Субстраты преимущественно кислые, с карбоксильной или другой ионизуемой группой. С 3A4 частично
перекрывается по липофильности, отсюда ранговая корреляция 0.69.

\emph{CYP1A2} --- узкая плоская щель. Субстраты --- плоские ароматические системы, полициклика
и азотистые гетероциклы. Они же его и индуцируют, из-за чего у курильщиков клиренс некоторых препаратов
заметно выше.

\emph{CYP2D6} объясняет половину этого документа. В активном центре сидят остатки аспартата и глутамата,
дающие солевой мостик с протонированным атомом азота лиганда. Практически все известные субстраты
и ингибиторы 2D6 --- основания, заряженные при физиологическом pH. Узнавание идёт не по общей
липофильности, а по конкретному электростатическому взаимодействию на конкретном расстоянии
от ароматического ядра.

Долю протонированной формы даёт уравнение Гендерсона--Хассельбальха:
\begin{equation}
f_{\text{прот}} = \frac{1}{1 + 10^{\,\mathrm{pH} - \mathrm{p}K_a}} .
\label{eq:hh}
\end{equation}
При pH 7.4 третичный амин с $\mathrm{p}K_a$ около 9.8 заряжен на 99.6 процента, пиридиновый азот
с $\mathrm{p}K_a$ 5.2 --- на 0.6 процента, а амид не заряжен вовсе. Разница в сто с лишним раз
по доле заряженных молекул. Фингерпринт, считающий фрагменты в радиусе двух связей, видит во всех трёх
случаях просто «атом азота в таком-то окружении» и различить их не может --- отсюда механистический
блок признаков.

И ещё: 2D6 сильно полиморфен. Часть популяции несёт
неработающие аллели, часть --- дупликации гена; это делает предсказание ингибирования 2D6 клинически
особенно ценным, но на наши данные не влияет, поскольку измерения делаются на рекомбинантном ферменте
одного варианта.
